%%%%%%%%%%%%%%%%%%%%%%%%%%%%%%%%%%%%%%%%%%%%%%%%%%%%%%%%%%%%%%%%%%%%%%%%%%%%%%%
\chapter{Semantics Utility Rules\label{chap:SemanticsUtilityRules}}
%%%%%%%%%%%%%%%%%%%%%%%%%%%%%%%%%%%%%%%%%%%%%%%%%%%%%%%%%%%%%%%%%%%%%%%%%%%%%%%

This chapter defines the following helper relations for operating on \nativevaluesterm{},
\hyperlink{def-envs}{environments}, and operations involving values and types:
\begin{itemize}
  \item \SemanticsRuleRef{GetPendingCalls} retrieves the value of the $\pendingcalls$ value from an environment;
  \item \SemanticsRuleRef{SetPendingCalls} updates the value of the $\pendingcalls$ value in an environment;
  \item \SemanticsRuleRef{IncrPendingCalls} updates the value of the $\pendingcalls$ value in an environment by incrementing it by one;
  \item \SemanticsRuleRef{DecrPendingCalls} updates the value of the $\pendingcalls$ value in an environment by decrementing it by one;
  \item \SemanticsRuleRef{RemoveLocal} removes a local storage element from a given environment;
  \item \SemanticsRuleRef{ReadIdentifier} creates an \executiongraphterm{} representing a Read Effect for a given identifier;
  \item \SemanticsRuleRef{WriteIdentifier} creates an \executiongraphterm{} representing a Write Effect for a given identifier;
  \item \SemanticsRuleRef{ConcatBitvectors} concatenates a list of bitvector \nativevaluesterm{};
  \item \SemanticsRuleRef{SlicesToPositions} converts a list of slices into a list of indices;
  \item \SemanticsRuleRef{MaxPosOfSlice} finds the maximum position in a slice;
  \item \SemanticsRuleRef{ReadFromBitvector} slices a bitvector \nativevalueterm{};
  \item \SemanticsRuleRef{WriteToBitvector} updates a slice of a bitvector \nativevalueterm{};
  \item \SemanticsRuleRef{GetIndex} reads a vector \nativevalueterm{} at a given index;
  \item \SemanticsRuleRef{SetIndex} updates a vector \nativevalueterm{} at a given index;
  \item \SemanticsRuleRef{GetField} reads a given field in a \nativevalueterm{} record;
  \item \SemanticsRuleRef{SetField} updates a given field in a \nativevalueterm{} record;
  \item \SemanticsRuleRef{DeclareLocalIdentifier} updates an environment by associating an identifier with a given \nativevalueterm{}
        and also generating an \executiongraphterm{} with a corresponding Write Effect;
  \item \SemanticsRuleRef{DeclareLocalIdentifierM} updates an environment by associating an identifier with a given \nativevalueterm{}
        and also updating a given \executiongraphterm{} with a corresponding Write Effect and an $\asldata$ edge;
  \item \SemanticsRuleRef{DeclareLocalIdentifierMM} updates an environment by associating an identifier with a given \nativevalueterm{}
        and also updating a given \executiongraphterm{} with a corresponding Write Effect and an $\aslpo$ edge.
\end{itemize}

\SemanticsRuleDef{GetPendingCalls}
\RenderRelation{get_pending_calls}

\ExampleDef{Retrieving the Number of Pending Calls}
In \listingref{CheckRecurseLimit}, the number of pending calls upon the call to \verb|factorial| from \verb|main|
with the environment $\env_0 \eqdef (\tenv_0, \denv_0)$ is $0$.\\
That is, $\getpendingcalls(\denv_0, \factorial) \evalarrow 0$.

\RenderProseAndFormally{get_pending_calls}

\SemanticsRuleDef{SetPendingCalls}
\RenderRelation{set_pending_calls}

\ExampleDef{Setting the Number of Pending Calls}
In \listingref{CheckRecurseLimit}, the number of pending calls upon the call to \verb|factorial| from \verb|main|
with the environment $\env_0 \eqdef (\tenv_0, \denv_0)$ is $0$, and it is then set to $1$.
That is, \\
$\setpendingcalls(\denv_0, \factorial, 1) \evalarrow \denv_1$
and\\
$\getpendingcalls(\denv_1, \factorial) \evalarrow 1$.

\RenderProseAndFormally{set_pending_calls}

\SemanticsRuleDef{IncrPendingCalls}
\RenderRelation{incr_pending_calls}

\ExampleDef{Incrementing the Number of Pending Calls}
In \listingref{CheckRecurseLimit}, the number of pending calls upon the call to \verb|factorial| from \verb|main|
with the environment $\env_0 \eqdef (\tenv_0, \denv_0)$ is $0$, and it is incremented to $1$.
That is, \\
$\incrpendingcalls(\denv_0, \factorial) \evalarrow \denv_1$
and\\
$\getpendingcalls(\denv_1, \factorial) \evalarrow 1$.

\RenderProseAndFormally{incr_pending_calls}

\SemanticsRuleDef{DecrPendingCalls}
\RenderRelation{decr_pending_calls}

\ExampleDef{Decrementing the Number of Pending Calls}
In \listingref{CheckRecurseLimit}, the number of pending calls upon returning from \verb|main| to \verb|factorial|
with the environment $\env_1 \eqdef (\tenv_1, \denv_1)$ is $1$, and it is decremented to $0$.
That is, \\
$\decrpendingcalls(\denv_1, \factorial) \evalarrow \denv_0$
and\\
$\getpendingcalls(\denv_0, \factorial) \evalarrow 0$.

\RenderProseAndFormally{decr_pending_calls}

\SemanticsRuleDef{RemoveLocal}
\RenderRelation{remove_local}

This relation is used to maintain the invariant that the dynamic environment
maintains bindings only for identifiers that are in scope. That is, only identifiers
that would be in the local dynamic environment.
%
This invariant is not necessary for correctness.
Removing identifiers from the dynamic environment does not have any
observable effect.

\ExampleDef{Removing Local Storage Elements from Dynamic Environments}
In \listingref{RemoveLocal}, the local storage element \verb|i| is removed after the (only) \forstatementterm,
and the local storage element \verb|exn| is removed after the (only) \trystatementterm.

\ASLListing{Removing local storage elements from dynamic environments}{RemoveLocal}{\semanticstests/SemanticsRule.RemoveLocal.asl}
The output to the console is the following:
% CONSOLE_BEGIN aslref \semanticstests/SemanticsRule.RemoveLocal.asl
\begin{Verbatim}[fontsize=\footnotesize, frame=single]
i=0
i=11/2
exn=TRUE
\end{Verbatim}
% CONSOLE_END

\RenderProseAndFormally{remove_local}

\SemanticsRuleDef{ReadIdentifier}
\RenderRelation{read_identifier}

\ExampleDef{The Effect of Reading from an Identifier}
In \listingref{ReadIdentifier}, the first iteration of the (only) \forstatementterm{}
generates (see \SemanticsRuleRef{EvalFor}) the following transition at the first iteration:
\[
\readidentifier(\vi, \nvint(0)) \evalarrow ( \overtext{\{\ReadEffect(\vi)\}}{\vnodes}, \overtext{\emptyset}{\vedges}, \overtext{\{\ReadEffect(\vi)\}}{\voutputnodes} ) \enspace.
\]

\ASLListing{A simple for loop}{ReadIdentifier}{\semanticstests/SemanticsRule.ReadIdentifier.asl}

\RenderProseAndFormally{read_identifier}

\SemanticsRuleDef{WriteIdentifier}
\RenderRelation{write_identifier}

\ExampleDef{The Effect of Writing to an Identifier}
In \listingref{ReadIdentifier}, the first iteration of the (only) \forstatementterm{}
generates (see \SemanticsRuleRef{EvalFor}) the following transition at the first iteration,
as the value of the index variable $\vi$ is incremented from $\nvint(0)$ to $\nvint(1)$:
\[
\writeidentifier(\vi, \nvint(1)) \evalarrow ( \overtext{\{\WriteEffect(\vi)\}}{\vnodes}, \overtext{\emptyset}{\vedges}, \overtext{\{\WriteEffect(\vi)\}}{\voutputnodes} ) \enspace.
\]

\RenderProseAndFormally{write_identifier}

\SemanticsRuleDef{ConcatBitvectors}
\RenderRelation{concat_bitvectors}

\ExampleDef{Concatenating Bitvectors}
The specification in \listingref{ConcatBitvectors} shows examples of concatenating
the \bitvectortypeterm{} fields of a \collectiontypeterm{} and of a \recordtypeterm.
\ASLListing{Concatenating bitvectors}{ConcatBitvectors}{\semanticstests/SemanticsRule.ConcatBitvectors.asl}

\RenderProseAndFormally{concat_bitvectors}

\SemanticsRuleDef{SlicesToPositions}
\RenderRelation{slices_to_positions}

\ExampleDef{Converting Slices to Lists of Indices}
In \listingref{SlicesToPositions}, the slices \verb|1+:3, 7:5|
are converted into the list of positions $3, 2, 1, 7, 6, 5$.
\ASLListing{Converting slices to lists of indices}{SlicesToPositions}{\semanticstests/SemanticsRule.SlicesToPositions.asl}

The specification in \listingref{SlicesToPositions-bad} terminates with a \dynamicerrorterm,
since the slice \verb|from+:6| evaluates to $(\nvint(-1), \nvint(6))$ and the starting position $-1$
is illegal for a slice.
\ASLListing{An illegal slice}{SlicesToPositions-bad}{\semanticstests/SemanticsRule.SlicesToPositions.bad.asl}

% Transliteration note: there are two slices_to_positions functions:
% the one in Native.ml, which calls the one in ASTUtils.ml.
% This rule transliterates the version in Native.ml, inlining the version
% in ASTUtils.ml.
\RenderProseAndFormally{slices_to_positions}

\SemanticsRuleDef{MaxPosOfSlice}
\RenderRelation{max_pos_of_slice}

\ExampleDef{Maximum positions of slices}
The slice \verb|1+:3| has maximum position 3.
The slice \verb|42+:0| has maximum position 42.

\RenderProseAndFormally{max_pos_of_slice}

\SemanticsRuleDef{ReadFromBitvector}
\RenderRelation{read_from_bitvector}

Notice that the bits of a bitvector go from the least-significant bit on the right
to the most-significant bit on the left.
This is reflected by how the rules list the bits.
The effect of placing the bits in sequence is that of concatenating the results
from all of the given slices.

\ExampleDef{Reading From a Bitvector}
The specification in \listingref{ReadFromBitvector} shows examples of reading slices from a bitvector
and reading slices from an integer (first converted to a bitvector), followed by the output to the console.

\ASLListing{Reading from a bitvector}{ReadFromBitvector}{\semanticstests/SemanticsRule.ReadFromBitvector.asl}
% CONSOLE_BEGIN aslref \semanticstests/SemanticsRule.ReadFromBitvector.asl
\begin{Verbatim}[fontsize=\footnotesize, frame=single]
empty_bv_slice = 0x, empty_i_slice = 0x
slice_bv = 0x9c5, slice_i = 0x9c5
\end{Verbatim}
% CONSOLE_END

\ExampleDef{Conversion of Integers to Bitvectors}
The specification in \listingref{AsBitvector} shows an example of converting a negative
integer into a bitvector, followed by the output to the console.

\ASLListing{Converting a signed integer into a bitvector}{AsBitvector}{\semanticstests/SemanticsRule.AsBitvector.asl}
% CONSOLE_BEGIN aslref \semanticstests/SemanticsRule.AsBitvector.asl
\begin{Verbatim}[fontsize=\footnotesize, frame=single]
bv = 0xfdf0, bv_i = 0xfdf0
\end{Verbatim}
% CONSOLE_END

\RenderProseAndFormally{read_from_bitvector}

\SemanticsRuleDef{WriteToBitvector}
\RenderRelation{write_to_bitvector}

See \ExampleRef{Writing to a Bitvector}, following the definition of $\writetobitvector$.

\RenderProseAndFormally{write_to_bitvector}

\ExampleDef{Writing to a Bitvector}
In reference to \listingref{semantics-leslice}, we have the following application of the current rule:
\begin{mathpar} % SUPPRESS_TEXTTT_LINTER
\inferrule{
  \src = \overname{0}{\vs_5}\overname{0}{\vs_4}\overname{0}{\vs_3}\overname{0}{\vs_2}\overname{0}{\vs_1}\overname{0}{\vs_0}\\
  \dst = \overname{1}{\vd_7}\overname{1}{\vd_6}\overname{1}{\vd_5}\overname{1}{\vd_4}\overname{1}{\vd_3}\overname{1}{\vd_2}\overname{1}{\vd_1}\overname{1}{\vd_0}\hva\\
  \slicestopositions(8, [\overname{(0, 4)}{\texttt{3:0}}, \overname{(6, 2)}{\texttt{7:6}}]) \evalarrow
  [3, 2, 1, 0, 7, 6]\hva\\
  \positions \eqdef [\overname{3}{I_5}, \overname{2}{I_4}, \overname{1}{I_3}, \overname{0}{I_2}, \overname{7}{I_1}, \overname{6}{I_0}]\\
  {\bitfunc = \lambda i \in 7..0.\left\{ \begin{array}{ll}
    \vs_j & \exists j\in 5..0.\ i = I_j\\
    \vd_i & \text{otherwise}
  \end{array} \right.}\hva\\
  \vbits \eqdef \bitfunc(7)\ \bitfunc(6)\ \bitfunc(5)\ \bitfunc(4)\ \bitfunc(3)\ \bitfunc(2)\ \bitfunc(1)\ \bitfunc(0)
}{
  {
  \begin{array}{r}
    \writetobitvector(
      [\overname{(0, 4)}{\texttt{3:0}}, \overname{(6, 2)}{\texttt{7:6}}],
      \overname{\nvbitvector(000000)}{\src},
      \overname{\nvbitvector(11111111)}{\dst}) \evalarrow\\
    \nvbitvector(
      \overname{0}{\vs_1}
      \overname{0}{\vs_0}
      \overname{1}{\vd_5}
      \overname{1}{\vd_4}
      \overname{0}{\vs_5}
      \overname{0}{\vs_4}
      \overname{0}{\vs_3}
      \overname{0}{\vs_2})
  \end{array}
  }
}
\end{mathpar}

\SemanticsRuleDef{GetIndex}
\RenderRelation{get_index}

\ExampleDef{Reading a Vector at an Index}
The specification in \listingref{GetIndex} shows examples of accessing indices of native vectors:
\begin{itemize}
  \item evaluating \verb|assert arr[[1]] == 3| involves \\
        $\getindex(1, \NVVector(\nvint(0), \nvint(3), \nvint(0))) \evalarrow \nvint(3)$;
  \item evaluating \verb|assert (5, 7).item0 == 5| involves \\
        $\getindex(0, \NVVector(\nvint(5), \nvint(7))) \evalarrow \nvint(5)$;
  \item evaluating \verb|(5, 7) as (integer, integer)| involves \\
        $\getindex(0, \NVVector(\nvint(5), \nvint(7))) \evalarrow \nvint(5)$ and\\
        $\getindex(1, \NVVector(\nvint(5), \nvint(7))) \evalarrow \nvint(7)$ to check the asserting type conversion.
\end{itemize}
\ASLListing{Reading a vector at an index}{GetIndex}{\semanticstests/SemanticsRule.GetIndex.asl}

\RenderProseAndFormally{get_index}

\SemanticsRuleDef{SetIndex}
\RenderRelation{set_index}

\ExampleDef{Writing a Vector at an Index}
In \listingref{GetIndex}, evaluating the statement \verb|arr[[1]] = 3;|
involves the following premise:
\[
\begin{array}{r}
\setindex(1, \nvint(3), \NVVector(\nvint(0), \nvint(0), \nvint(0))) \evalarrow\\
\NVVector(\nvint(0), \nvint(3), \nvint(0)) \enspace.
\end{array}
\]

\RenderProseAndFormally{set_index}

\SemanticsRuleDef{GetField}
\RenderRelation{get_field}

\ExampleDef{Reading a Record Field}
In \listingref{GetField}, there are the following examples of reading a record field:
\begin{itemize}
  \item evaluating the expression \verb|color_to_int[[RED]]| involves\\
        $\getfield(\RED, \NVRecord(\RED\mapsto \nvint(0), \GREEN\mapsto \nvint(1), \BLUE\mapsto \nvint(2))) \evalarrow \nvint(0)$;
  \item evaluating the expression \verb|r.RED| involves\\
        $\getfield(\RED, \NVRecord(\RED\mapsto \nvint(0), \GREEN\mapsto \nvint(1), \BLUE\mapsto \nvint(2))) \evalarrow \nvint(0)$.
\end{itemize}
\ASLListing{Reading a record field}{GetField}{\semanticstests/SemanticsRule.GetField.asl}

\RenderProseAndFormally{get_field}
The typechecker ensures, via \TypingRuleRef{EGetRecordField}, that the field $\name$ exists in $\record$.

\SemanticsRuleDef{SetField}
\RenderRelation{set_field}

\ExampleDef{Writing to a Record Field}
In \listingref{GetField}, there are the following examples of writing to a record field:
\begin{itemize}
  \item evaluating the statement \verb|color_to_int[[GREEN]] = 1;| involves\\
        $
        \begin{array}{r}
        \setfield\left(
          \begin{array}{l}
          \GREEN, \\
          \nvint(1),\\
          \NVRecord(\RED\mapsto \nvint(0), \GREEN\mapsto \nvint(0), \BLUE\mapsto \nvint(0))
          \end{array}
          \right) \evalarrow\\
        \NVRecord(\RED\mapsto \nvint(0), \GREEN\mapsto \nvint(1), \BLUE\mapsto \nvint(0)) \enspace;
        \end{array}
        $
  \item evaluating the statement \verb|r.BLUE = -1;| involves\\
        $
        \begin{array}{r}
        \setfield\left(
          \begin{array}{l}
          \BLUE, \\
          \nvint(-1),\\
          \NVRecord(\RED\mapsto \nvint(0), \GREEN\mapsto \nvint(1), \BLUE\mapsto \nvint(2))
          \end{array}
          \right) \evalarrow\\
        \NVRecord(\RED\mapsto \nvint(0), \GREEN\mapsto \nvint(1), \BLUE\mapsto \nvint(-1)) \enspace.
        \end{array}
        $
\end{itemize}

The typechecker ensures that the field $\name$ exists in $\record$.

\RenderProseAndFormally{set_field}

\SemanticsRuleDef{DeclareLocalIdentifier}
\RenderRelation{declare_local_identifier}

\ExampleDef{Evaluating Local Declarations}
In \listingref{semantics-levar}, evaluating \verb|var x: integer = 3;|
binds \verb|x| to $\nvint(3)$ in the environment where \verb|x| is unbound
as well as producing the \executiongraphterm{} $\WriteEffect(\vx)$.

\RenderProseAndFormally{declare_local_identifier}

\SemanticsRuleDef{DeclareLocalIdentifierM}
\RenderRelation{declare_local_identifier_m}

See \ExampleRef{Evaluating Local Declarations}.

\RenderProseAndFormally{declare_local_identifier_m}

\SemanticsRuleDef{DeclareLocalIdentifierMM}
\RenderRelation{declare_local_identifier_mm}

See \ExampleRef{Evaluating Local Declarations}.

\RenderProseAndFormally{declare_local_identifier_mm}
